% 355_SSB_Torus_De_ch.tex
% Johann Pascher, 7. September 2026
% Spontane Symmetriebrechung SU(2)_L x U(1)_Y -> U(1)_EM aus T4/Z3-Geometrie
\providecommand{\BM}{\textbf{[B]}}
\providecommand{\KM}{\textbf{[K]}}
\providecommand{\SM}{\textbf{[S]}}
\providecommand{\XM}{\textbf{[X]}}
\providecommand{\QM}{\textbf{[Q]}}

\phantomsection
\addcontentsline{toc}{chapter}{%
  SSB $SU(2)_L\times U(1)_Y\to U(1)_\text{EM}$ aus $T^4/Z_3$-Geometrie}

\section*{Statuszeichen}
\begin{center}
\begin{tabular}{@{}lp{10.5cm}@{}}
\textbf{[K]} & \emph{Konfirmiert} --- numerisch geprüft. \\[3pt]
\textbf{[B]} & \emph{Belegt} --- algebraisch bewiesen. \\[3pt]
\textbf{[S]} & \emph{Spekulativ} --- offen. \\[3pt]
\textbf{[X]} & \emph{Experimentell} --- gesichert. \\[3pt]
\textbf{[Q]} & \emph{Quelle} --- externer Primärwert. \\
\end{tabular}
\end{center}

\section*{Abstract}

Die spontane Symmetriebrechung $SU(2)_L\times U(1)_Y\to U(1)_\text{EM}$
ist im Standardmodell durch das Higgs-Potential $V(\Phi)$ mit $\mu^2<0$
implementiert --- eine Setzung. Dieses Dokument leitet sie aus drei
Bausteinen der $T^4/Z_3$-Geometrie ab, die alle bereits im Korpus belegt sind.

\textbf{Die drei Bausteine:}
\begin{enumerate}
\item $\tilde{T}\cdot m=1$ erzwingt $\langle\Phi\rangle\neq0$: das Vakuum
  muss Masse tragen \BM\ (Dok.~003, 010, 306, 334).
\item Die $Z_3$-Galois-Struktur erzwingt $Q_\text{vak}=0$: das Vakuum
  ist elektrisch neutral \BM\ (Dok.~346).
\item $Q=0$ des Vakuums selektiert $U(1)_\text{EM}$ als Restgruppe:
  der Gell-Mann-Nishijima-Generator annihiliert das Vakuum \BM\ (Dok.~347).
\end{enumerate}

\textbf{Hauptbefund:} Die Symmetriebrechungsstruktur ist topologisch erzwungen,
nicht postuliert \BM. Was nach SM-Konvention als spontan brechende Setzung
$\mu^2<0$ erscheint, ist in FFGFT eine Konsequenz der Zeitfeld-Geometrie und
der Galois-Klassifikation.

\textbf{Offen bleibt \SM:} Die CKM/PMNS-Mischungswinkelwerte (Dok.~348 Satz~D).

\section{Ausgangslage}

\subsection*{Das SM-Higgs-Potential}

Das Standardmodell implementiert SSB durch:
\begin{equation}
  V(\Phi) = \mu^2\,|\Phi|^2 + \lambda\,|\Phi|^4,
  \quad \mu^2 < 0,
  \label{eq:higgs_pot}
\end{equation}
mit dem Minimum $|\Phi|^2 = -\mu^2/(2\lambda) = v^2/2$, $v\approx246\,\text{GeV}$.
Das Higgs-Dublett $\Phi = (\phi^+,\phi^0)^T$ hat $I_3 = (-\tfrac{1}{2},+\tfrac{1}{2})$,
$Y = +1$. Das neutrale Komponent $\phi^0$ entwickelt den VEV:
\begin{equation}
  \langle\Phi\rangle = \begin{pmatrix}0\\v/\sqrt{2}\end{pmatrix},
  \quad Q_\text{vak} = I_3 + \frac{Y}{2} = -\tfrac{1}{2}+\tfrac{1}{2} = 0.
  \label{eq:sm_vev}
\end{equation}
Die überlebende Symmetrie ist $U(1)_Q$ mit Generator $Q=I_3+Y/2$ ---
das ist $U(1)_\text{EM}$ \XM\ (PDG~\QM).

\textbf{Die Frage:} Warum hat das Vakuum $Q=0$? Warum $\mu^2<0$?
Im SM: Setzung. In FFGFT: aus der Geometrie.

\section{Baustein 1: VEV $\neq 0$ aus $\tilde{T}\cdot m = 1$}

Die Zeitfeld-Dualität (Dok.~003, 010) lautet in natürlichen Einheiten:
\begin{equation}
  \tilde{T}(x,t)\cdot m(x,t) = 1.
  \label{eq:tm1}
\end{equation}
Diese Relation hat bei $m=0$ keine Lösung: $\tilde{T}\to\infty$ ist
kein physikalischer Zustand (Dok.~306, 334). Das Vakuum muss eine
nichtverschwindende Masse tragen:
\begin{equation}
  m_\text{vak} \neq 0 \quad\Longleftrightarrow\quad \langle\Phi\rangle \neq 0.
  \label{eq:vev_nonzero}
\end{equation}

\textbf{Verbindung zum Higgs-Potential:} In SM-Sprache ist $m_\text{vak}=0$ das
Minimum von $V$ für $\mu^2>0$ --- ein leeres Vakuum. Für $\mu^2<0$ liegt das
Minimum bei $|\Phi|\neq0$. Die Bedingung \eqref{eq:vev_nonzero} aus FFGFT ist
äquivalent zur Forderung $\mu^2<0$ im SM-Potential: beide erzwingen
$\langle\Phi\rangle\neq0$ \BM.

\textbf{Der Unterschied:} Im SM ist $\mu^2<0$ eine Setzung.
In FFGFT folgt $m_\text{vak}\neq0$ aus der Reziprozität $\tilde{T}\cdot m=1$,
die selbst aus $E=mc^2$ folgt (Dok.~354 Erkenntnis~1) \BM.

\section{Baustein 2: $Q_\text{vak} = 0$ aus der Galois-Struktur}

\subsection*{Schritt 2a: Das Vakuum ist ein Torus-Hintergrund}

Das Zeitfeld im Vakuum ist $\tilde{T}_0 = 1/m_0 = \text{const}$ --- eine
statische, räumlich homogene Konfiguration. Im $T^4/Z_3$-Rahmen ist das die
$k=0$-Mode (Nullwindung): kein Drehimpuls, keine Farbladung, keine
elektrische Ladung. Das Vakuum ist topologisch trivial \BM\ (Dok.~A075).

\subsection*{Schritt 2b: Die Galois-Klassifikation des Vakuums}

In GF$(27)^*$ klassifiziert die Kombination aus QR-Eigenschaft mod~13
und $N\bmod3$-Farbladung alle Fermionen (Dok.~346 Satz~B \BM):

\begin{center}
\begin{tabular}{@{}llll@{}}
\toprule
QR mod~13 & $N\bmod3$ & Teilchentyp & $Q$ \\
\midrule
QR (Quadratischer Rest) & $\equiv 0$ & Neutrino & $0$ \\
NQR (Nicht-QR) & $\equiv 0$ & Geladenes Lepton & $-1$ \\
NQR & $\equiv 1,2$ & Quark & $\pm1/3,\pm2/3$ \\
\bottomrule
\end{tabular}
\end{center}

Das Vakuum ist kein Fermion --- es ist die $k=0$-Mode, das topologisch triviale
Element. In GF$(27)^*$ entspricht das dem Einselement (multiplikative Identität),
das zu jedem Orbit $O_i$ gehört als Grenzfall. Die entscheidende Eigenschaft:

\begin{itemize}
  \item Keine Farbladung: $N\bmod3 = 0$ (Farb-Singlett) \BM
  \item Elektrisch neutral: QR-Eigenschaft, also $Q=0$ \BM
\end{itemize}

Dies folgt aus dem Satz~A von Dok.~346 \BM:
$Q=0 \iff \text{QR mod~13}$.
Das topologisch triviale Element (Nullwindung) ist QR \BM.

\subsection*{Schritt 2c: Das Vakuum hat $I_3=-1/2$, $Y=+1$ mit $Q=0$}

Im SM hat das Higgs-Dublett $Y=+1$, und das neutrale Komponent
$\phi^0$ hat $I_3=-1/2$. Dann ist $Q=I_3+Y/2=0$.

In FFGFT ergibt sich die Dublett-Struktur aus dem Sd-Sektor (Dok.~349 Satz~D \BM):
Der Sd-Sektor (linkshändig, ungerade Cl$(6)$-Grade) trägt $I_3=-1/2$.
Die Hyperladung $Y$ folgt aus dem Legendre-Symbol (Dok.~347 Satz~A \BM):
$Y = -1 + \tfrac{4}{3}\,\text{NQR}$; für das Vakuum mit $\text{QR}$
(elektrisch neutral) ist $Y=+1$ wenn der Sektor die volle Hyperladung trägt.

Das Vakuum liegt im Sd-Sektor mit $I_3=-1/2$, $Y=+1$:
\begin{equation}
  Q_\text{vak} = I_3 + \frac{Y}{2} = -\frac{1}{2} + \frac{1}{2} = 0.
  \label{eq:qvak}
\end{equation}
$Q_\text{vak}=0$ ist algebraisch erzwungen \BM.

\section{Baustein 3: $U(1)_\text{EM}$ als Restgruppe}

\subsection*{Residualsymmetrie des Vakuums}

Die Symmetriegruppe $G = SU(2)_L\times U(1)_Y$ wirkt auf das Vakuum.
Die Restgruppe $H\subset G$ nach der Brechung ist die Untergruppe, die
das Vakuum invariant lässt:
\begin{equation}
  H = \{g\in G \mid g\,|\text{vak}\rangle = |\text{vak}\rangle\}.
  \label{eq:residual}
\end{equation}

Der Generator $Q = I_3 + Y/2$ (Gell-Mann-Nishijima, Dok.~347 Satz~C \BM)
annihiliert das Vakuum:
\begin{equation}
  Q\,|\text{vak}\rangle = Q_\text{vak}\,|\text{vak}\rangle = 0.
  \label{eq:Q_annihilates}
\end{equation}
Die zugehörige Ein-Parameter-Untergruppe:
\begin{equation}
  U(1)_Q : \quad e^{i\alpha Q}\,|\text{vak}\rangle
  = e^{i\alpha\cdot 0}\,|\text{vak}\rangle = |\text{vak}\rangle
  \quad\forall\,\alpha\in\mathbb{R}.
  \label{eq:u1_invariant}
\end{equation}
$U(1)_Q$ lässt das Vakuum invariant --- das ist $U(1)_\text{EM}$ \BM.

\subsection*{Alle anderen Generatoren brechen die Symmetrie}

Die übrigen Generatoren von $SU(2)_L\times U(1)_Y$ --- nämlich $T_1$, $T_2$
(ladungsändernde Operatoren) und die Kombination $T_3-\frac{Y}{2}$ ---
haben nichttriviale Eigenwerte auf dem Vakuum und erzeugen die massiven
Bosonen $W^\pm$, $Z^0$ (Goldstone-Theorem \XM\ \QM):
\begin{align}
  W^\pm &\leftrightarrow T_\pm = T_1 \pm iT_2, \quad
    T_\pm|\text{vak}\rangle \neq 0 \label{eq:W}\\
  Z^0 &\leftrightarrow T_3 - \frac{Y}{2}\sin^2\theta_W, \quad
    Z^0|\text{vak}\rangle \neq 0. \label{eq:Z}
\end{align}

Das Photon $\gamma$ entspricht dem Generator $Q = T_3 + Y/2$, der das
Vakuum annihiliert --- es bleibt masselos \BM\ (Goldstone-Theorem).
Die Massenverhältnisse $M_W/M_Z = \cos\theta_W$ folgen aus dem
Weinberg-Winkel, der aus $\xi$ hergeleitet ist (Dok.~323 \KM).

\section{Die vollständige Kette}

\begin{center}
\begin{tabular}{@{}p{0.6cm}p{7.5cm}p{1.5cm}p{3.5cm}@{}}
\toprule
Schritt & Aussage & Status & Herkunft \\
\midrule
1 & $\tilde{T}\cdot m=1$ aus $E=mc^2$ & \BM & Dok.~003,010,052 \\
2 & $m_\text{vak}\neq0$: kein masseloses Vakuum & \BM & Dok.~306,334 \\
3 & Vakuum ist $k=0$-Mode (topologisch trivial) & \BM & Dok.~A075 \\
4 & $Q_\text{vak}=0$: topologisch triviales Element ist QR & \BM & Dok.~346 Satz A \\
5 & Sd-Sektor: $I_3=-1/2$ (linkshändig) & \BM & Dok.~349 Satz D \\
6 & $Y=+1$ aus Legendre-Symbol für neutrales Vakuum & \BM & Dok.~347 Satz A \\
7 & $Q=I_3+Y/2=0$ (Gell-Mann-Nishijima) & \BM & Dok.~347 Satz C \\
8 & $U(1)_Q$ lässt Vakuum invariant: Restgruppe & \BM & dieses Dok. \\
9 & $W^\pm$, $Z^0$ massiv; $\gamma$ masselos & \BM+\XM & Goldstone + Dok.~323 \\
\midrule
\multicolumn{2}{@{}l}{SSB $SU(2)_L\times U(1)_Y\to U(1)_\text{EM}$} & \BM & vollständig \\
\bottomrule
\end{tabular}
\end{center}

\section{Numerische Konsistenz}

\begin{center}
\begin{tabular}{@{}lllll@{}}
\toprule
Größe & FFGFT & Experiment \QM & Abw. & Status \\
\midrule
$\sin^2\theta_W(M_Z)$ & $0{,}2308$ & $0{,}2312$ & $-0{,}17\,\%$ & \KM \\
$M_W/M_Z$ & $\cos\theta_W = 0{,}8813$ & $0{,}8819$ & $-0{,}07\,\%$ & \KM \\
$Q_\text{vak}$ & $0$ (exakt) & $0$ (exakt) & $0$ & \BM \\
$m_\gamma$ & $0$ (exakt) & $<10^{-18}$\,eV & --- & \BM+\XM \\
\bottomrule
\end{tabular}
\end{center}

\section{Was offen bleibt}

\textbf{CKM/PMNS-Winkelwerte \SM:}
Die Struktur der Mischungsmatrizen (welche Einträge null, welche $1/\sqrt{2}$)
ist belegt \BM\ (Dok.~348 Sätze B,C). Die Zahlenwerte der Mischungswinkel
erfordern weitere algebraische Arbeit (Dok.~348 Satz~D). Das ist das einzige
verbleibende \SM\ im Eich/Higgs-Komplex.

\textbf{Neuer Ansatz: GF$(3^6)$ = GF(729) und 7.\ Einheitswurzeln \SM:}

IPI-Korrespondenz mit Douglas Matzke (7.~September 2026) liefert einen
neuen algebraischen Befund: 7.\ Einheitswurzeln benötigen eine Erweiterung
auf GF$(3^6)$ = GF(729) \BM. Der Grund ist algebraisch zwingend: das Polynom
$x^7 - 1$ faktorisiert über GF$(3)$ als $(x-1)\cdot\phi_7$, wobei $\phi_7$
\emph{irreduzibel vom Grad~6} ist --- eine Grad-6-Erweiterung ist daher
minimal und notwendig.

GF$(729)^*$ hat Ordnung $728 = 8 \times 7 \times 13$. Die Primfaktoren
liefern drei Skalen:
\begin{itemize}
  \item $13$: GF$(27)$-Ebene --- Quantenzahlen, Ladungen, Generationen \BM
  \item $7$: heptagonale Struktur, neu; Phasen $k\cdot360°/7$ \BM
  \item $8$: oktagonale Struktur; $360°/8 = 45°$ --- nahe an
    $\theta_{23}^\text{PMNS}\approx49{,}2°$ (Kandidat \SM)
\end{itemize}

\textbf{Konkrete Beobachtung \SM:}
$\theta_{23}^\text{PMNS}\approx49{,}2°$ liegt $2{,}2°$ von $360°/7\approx51{,}4°$
entfernt --- innerhalb der Messunschärfe, aber kein Beweis.
Die CP-Phase $\delta_{CP}^\text{CKM}\approx68°$ hat keinen direkten Bezug zu
7.\ Einheitswurzeln; $\delta_{CP}^\text{PMNS}\approx195°\approx4\cdot360°/7-10°$
ist ebenfalls kein direkter Treffer.

\textbf{Was das bedeutet:} GF$(729)$ ist der natürliche nächste Schritt
über GF$(27)$. Falls CKM/PMNS-Winkel auf dieser Ebene liegen, wäre
Dok.~348 Satz~D auf GF$(729)$ auszuführen --- eine konkrete offene Aufgabe,
kein strukturelles Hindernis.

\textbf{Warum $P_+$ und nicht $P_-$ \SM:}
Die chirale Projektion $P_+$ selektiert den linkshändigen Sd-Sektor (A095).
Warum die Natur $P_+$ wählt (und nicht $P_-$), ist geometrisch plausibel
aber nicht aus einer tieferen Symmetriebrechung abgeleitet \SM.

\section{Prüfskript}

\begin{itemize}[leftmargin=2em,itemsep=2pt,topsep=2pt]
  \item Numerische Verifikation der SSB-Kette: $Q_\text{vak}=0$,
    Residualsymmetrie $U(1)_Q$, $W^\pm/Z^0$-Massenverhältnis aus $\theta_W$,
    Goldstone-Zählung:
    \texttt{python/Dok355\_Skripte/pruef\_355\_ssb.py}
\end{itemize}

\section{Verwendung von KI-Werkzeugen}

Dieses Dokument ist unter Verwendung KI-gestützter Werkzeuge entstanden.
Die Fragestellungen, die Setzungen und alle inhaltlichen Entscheidungen
stammen vom Autor; die Werkzeuge formulieren aus, rechnen nach und erstellen
Prüfskripte. Die Verantwortung für Inhalt und Fehler liegt vollständig beim
Autor. Der ausführliche Wortlaut steht in A010.
